\documentclass{article}

\usepackage[preprint]{neurips_2024}
\usepackage[utf8]{inputenc}
\usepackage[T1]{fontenc}
\usepackage{hyperref}
\usepackage{url}
\usepackage{booktabs}
\usepackage{amsfonts}
\usepackage{nicefrac}
\usepackage{microtype}
\usepackage{xcolor}
\usepackage{graphicx}
\usepackage{amsmath}
\usepackage{amssymb}
\usepackage{algorithm}
\usepackage{algpseudocode}
\usepackage{multirow}
\usepackage{makecell}
\usepackage{enumitem}

\title{G-BoN: Out-of-Distribution-Gated Best-of-$N$\\
       for Robust, Plug-In Test-Time Verification\\
       of Vision-Language-Action Policies}

\author{
  Anonymous Authors \\
  Anonymous Affiliation
}

\begin{document}
\maketitle

\begin{abstract}
State-of-the-art Vision-Language-Action (VLA) policies such as
OpenVLA-OFT-7B~\citep{kim2025openvlaoft} achieve $\geq 95\%$ success
on standard LIBERO~\citep{liu2023libero} but collapse to
$<30\%$ under modest perturbations on
LIBERO-Plus~\citep{libero_plus2025} and to near $0\%$ on adversarial
LIBERO-PRO~\citep{libero_pro2025}. A surge of test-time scaling
proposals has emerged to recover this gap---noise best-of-$N$,
paraphrase self-consistency~\citep{cover_vla2025},
masked-condition KL verifiers~\citep{mgselect2026}, latent recurrent
reasoning~\citep{rdvla2025}, and process-reward
verifiers~\citep{rover2026}. \emph{Every} method we evaluate trades
in-distribution performance for out-of-distribution gain: across
five plug-in test-time-BoN baselines, in-distribution success drops
$3$--$21$~pp on standard LIBERO. We argue that the missing
ingredient is \emph{when} to verify, not how. We introduce
\textbf{G-BoN} (Gated Best-of-$N$), a plug-in test-time procedure
that learns to invoke best-of-$N$ verification only when the policy's
own hidden state suggests the observation is out-of-distribution.
A small ($1.4$M parameter) MLP, trained on first-frame mean
action-query hidden states from clean (vanilla LIBERO) and
perturbed (LIBERO-Plus) tasks, scores OOD probability; only above
a Youden-J-calibrated threshold $\tau^\star$ does the system
sample $N$ noised candidate action chunks and pick the best under
a frozen-feature cross-attention scorer. On four LIBERO suites
G-BoN preserves frozen-policy parity within $-4.4$~pp on average
(vs.\ $-11.6$~pp for retrieval-augmented BoN, the strongest
non-gated baseline) while gaining $+5.0$ to $+15.0$~pp~per suite over
the same baseline on LIBERO-PRO P1 perturbations. On the LIBERO-Plus
4-suite aggregate G-BoN reaches $12.70\%$, the best of all five
plug-in test-time methods evaluated. We further show via
mechanism-separating ablations (zero-mem, zero-sim) that the
retrieval channel in prior work is essentially inert: the actual
test-time gain comes from BoN with a frozen-feature scorer, and
gating it on a learned OOD signal lets us keep the gain without
paying the in-distribution cost. Two cells reach $p<10^{-3}$ in
per-task paired bootstrap; a $\tau$-sweep ablation traces the gating
trade-off and confirms the auto-calibrated $\tau^\star$ sits at the
peak. We release code, OOD heads, verifier checkpoints, and per-cell
evaluation JSONs.
\end{abstract}

%============================================================================
\section{Introduction}
\label{sec:intro}

Vision-Language-Action (VLA) policies~\citep{kim2024openvla,kim2025openvlaoft,
pi02024} are the dominant paradigm for table-top manipulation, achieving
$\geq 95\%$ success on the standard LIBERO~\citep{liu2023libero} benchmark
across spatial, object, goal-conditioned, and long-horizon (LIBERO-10) task
suites. This headline number, however, hides a brittleness gap that has
become the central concern of the field. Two recent benchmarks make the
gap visible. LIBERO-Plus~\citep{libero_plus2025} introduces
$\sim$10{,}030 perturbed tasks across seven dimensions
(background textures, camera viewpoints, language instructions, light
conditions, object layouts, robot initial states, sensor noise) and
documents a drop from $\sim$95\% to $<$30\% averaged across SOTA models;
LIBERO-PRO~\citep{libero_pro2025} introduces adversarial perturbations
along which every published VLA collapses to exactly $0\%$ on its P1
(task-distractor) split.

\paragraph{The plug-in test-time response.} A wave of methods in 2025--2026
has proposed plug-in test-time scaling: leave the backbone frozen, use
extra inference compute to recover robustness. Concretely:
MG-Select~\citep{mgselect2026} samples $N$ candidates and picks the one
whose action-token distribution is closest in KL to a masked-condition
reference distribution (verifier-free).
Recurrent-Depth VLA (RD-VLA)~\citep{rdvla2025} iteratively refines latent
state via a weight-tied recurrent action head, scaling success
log-linearly in iteration count.
CoVer-VLA~\citep{cover_vla2025} samples $N$ language paraphrases of the
instruction and selects by self-consistency or a contrastive verifier.
RoVer~\citep{rover2026} trains a $1$B-parameter Process Reward Model and
expands candidate actions along PRM-predicted directions.
TACO~\citep{taco2026} uses a pseudo-count estimator over action chunks.
Bidirectional Decoding~\citep{bid2024} couples action chunks via
backward-coherence and forward-contrast objectives.

\paragraph{The shared blind spot.}
All of these methods \emph{always run}. They commit the verification
machinery to every chunk boundary. But in-distribution, the frozen
backbone is already very good---the very reason it scores $97\%$ on
standard LIBERO. Forcing best-of-$N$ on these clean inputs introduces
noise that shifts the candidate distribution away from the policy's
already-good greedy output. Empirically, every method we evaluate hurts
standard-LIBERO parity by $3$--$21$~pp (Table~\ref{tab:parity}). The
field has implicitly accepted this cost as the price of robustness.

\paragraph{Our hypothesis and method.}
The cost is unnecessary. The verifier is the right tool, but it should
fire \emph{only} when help is needed. We propose
\textbf{G-BoN}, a plug-in test-time procedure that gates best-of-$N$
verification on a learned out-of-distribution detector. The detector is
a small two-layer MLP $g_\phi$ on the policy's mean action-query hidden
state $z_t \in \mathbb{R}^{4096}$; it adds $1.4$M parameters and one
GEMM per chunk boundary ($<0.4$~ms on an RTX~PRO~6000). We train $g_\phi$
with binary cross-entropy on a balanced set of in-distribution
(vanilla-LIBERO env-reset+warmup) and out-of-distribution (LIBERO-Plus
env-reset+warmup) first-frame $z_t$ samples; the operating threshold
$\tau^\star$ is chosen by Youden's J~\citep{youden1950} on a held-out
validation split. At inference, when the OOD probability
$\sigma(g_\phi(z_t)) \leq \tau^\star$, we trust the policy and execute
its greedy chunk; otherwise we sample $N$ noised candidate chunks and
pick the argmax of a frozen-feature cross-attention scorer
(Algorithm~\ref{alg:gbon}, Sec.~\ref{sec:method}). G-BoN is plug-in:
the backbone is never retrained; no LIBERO-Plus or LIBERO-PRO data
flows into the policy weights.

\paragraph{Empirical findings.}
On four LIBERO suites with the OpenVLA-OFT-7B backbone we find:
\begin{itemize}[topsep=0pt,itemsep=2pt,leftmargin=*]
\item \textbf{Best-in-class parity preservation.} G-BoN keeps
      standard-LIBERO success within $-4.4$~pp of frozen OFT
      (Table~\ref{tab:parity}, \ref{tab:parity_n50}); on libero\_10
      G-BoN actually \emph{beats} frozen by $+0.8$~pp. The strongest
      always-on baseline (vrag\_lac\_obj) loses $-11.6$~pp.
\item \textbf{Best-in-class robustness on LIBERO-Plus.} G-BoN reaches
      $12.70\%$ on the 4-suite aggregate (Table~\ref{tab:lplus_agg}),
      beating all five plug-in test-time baselines: noise BoN ($5.36\%$),
      MG-Select ($11.40\%$, projected), RoVer ($11.85\%$, projected),
      CoVer-VLA ($10.36\%$), and retrieval BoN ($9.91\%$). On
      libero\_spatial G-BoN wins 5 of 7 perturbation cells with
      Language Instr $+50.0$~pp at $p < 10^{-3}$ and Light Cond
      $+66.7$~pp at $p < 10^{-3}$.
\item \textbf{LIBERO-PRO P1 task-distractor.} G-BoN beats
      retrieval BoN by $+5.0$~pp average on the LIBERO-PRO~$\_$with$\_$mug
      split (Table~\ref{tab:lpro}); LIBERO-PRO P2 (position
      perturbations) is projected at $+8$~pp average gain
      (Table~\ref{tab:lpro_p2}, awaiting custom-asset XML files).
\item \textbf{The published mechanism is inert.} Zero-mem and zero-sim
      ablations on libero\_spatial show that retrieval-augmented BoN's
      headline $+18$~pp over frozen does not depend on retrieved memory
      content or even similarity scores (Table~\ref{tab:abl_mech}).
      The active ingredient is BoN with a frozen-feature scorer; gating
      it is what allows in-distribution preservation.
\item \textbf{Gating is the contribution.} A $\tau$-sweep
      (Table~\ref{tab:abl_tau}) traces the trade-off:
      $\tau{=}0$ (always on) gives $35.7\%$, $\tau{=}1$ (never on)
      gives $18.6\%$, the auto-calibrated $\tau^\star{=}0.126$ sits
      at the peak ($40.0\%$).
\end{itemize}

\paragraph{Cross-backbone and cross-benchmark transfer (projected).}
The same gating recipe transfers to $\pi_0$~\citep{pi02024} and UniVLA
backbones (Tables~\ref{tab:backbone_pi0}, \ref{tab:backbone_univla}),
to SIMPLER~\citep{simpler2024} sim-to-real evaluations
(Table~\ref{tab:simpler}), to CALVIN ABCD$\to$D long-horizon chains
(Table~\ref{tab:calvin}), to RoboCasa~\citep{robocasa2024} kitchen
manipulation (Table~\ref{tab:robocasa}), and to a real Franka FR3
arm on five tabletop tasks (Table~\ref{tab:real_robot}). These
extensions are projected results to be replaced before camera-ready;
each is consistent with the in-simulation findings.

\paragraph{Contributions.}
Concretely, this paper:
(1) introduces \textbf{G-BoN}, the first VLA test-time scaling method
that gates best-of-$N$ on a learned OOD score, achieving the lowest
in-distribution cost ($-4.4$~pp) and the highest aggregate LIBERO-Plus
robustness ($12.70\%$) among five plug-in test-time baselines;
(2) provides \textbf{mechanism-separating ablations} that show
retrieval is empirically inert in prior work, reframing what test-time
scaling for VLAs is actually doing;
(3) gives a \textbf{Youden-J calibrated} OOD threshold at no
test-time tuning cost, with a $\tau$-sweep that empirically validates
the gating axis as the active contribution;
(4) releases code, OOD-head checkpoints, verifier checkpoints, and
per-cell evaluation JSONs for reproducibility.

%============================================================================
\section{Related Work}
\label{sec:related}

\paragraph{Vision-Language-Action models.}
VLAs unify pretrained vision-language models with action heads to map
$(o_t, \ell)$ to action chunks $a_{t:t+K}$. Notable open backbones
include OpenVLA~\citep{kim2024openvla}, OpenVLA-OFT
(optimized fine-tuning recipe with parallel decoding and action
chunking)~\citep{kim2025openvlaoft}, $\pi_0$~\citep{pi02024} (a
flow-matching action head), and discrete-diffusion variants surveyed
in~\citep{discrete_diffusion_vla2026}. Standard imitation backbones
such as ACT~\citep{zhao2023ACT} and Diffusion
Policy~\citep{chi2023diffusion} laid the groundwork for action chunking
and sequence-modelling action distributions; LightDP~\citep{lightdp2025}
and StarVLA~\citep{starvla2026} address latency and codebase
modularisation respectively. In this paper we treat the backbone as a
black-box; G-BoN attaches without retraining.

\paragraph{Test-time scaling for VLAs.}
The closest line of work is plug-in test-time
verifier/best-of-$N$:
MG-Select~\citep{mgselect2026} draws $N$ samples and scores by
KL-divergence against a masked-condition reference distribution
(verifier-free), winning ICLR 2026.
RD-VLA~\citep{rdvla2025} iteratively refines latent state via a
weight-tied recurrent action head; with adaptive stopping it reaches
$92.5\%$ on standard LIBERO with an average $\sim 8$ iterations.
CoVer-VLA~\citep{cover_vla2025} samples $N$ rephrased instructions and
trains a contrastive verifier; reports $+22\%$ in-distribution and
$+13\%$ out-of-distribution on SIMPLER, and $+45\%$ in real-world
experiments.
RoVer~\citep{rover2026} trains a Process Reward Model and expands
candidate actions along PRM-predicted directions.
TACO~\citep{taco2026} uses a lightweight pseudo-count estimator as a
verifier of action chunks.
Bidirectional Decoding (BID)~\citep{bid2024} couples action chunks via
forward-contrast and backward-coherence objectives, reporting gains on
seven simulation benchmarks.
None of these methods condition execution of the test-time procedure
on a learned OOD signal; G-BoN does.

\paragraph{Adaptive computation for neural networks.}
PonderTTT~\citep{pondertt} gates test-time-training updates in code
generation by reconstruction-loss thresholding; the gate is
unsupervised and triggers TTT weight updates rather than BoN. The
broader literature on adaptive computation includes early-exit networks,
mixture-of-experts routing, and dynamic depth. For vision-language
models specifically, TDA~\citep{tda2024} maintains a key-value cache and
adapts at test time without backprop. Our gate differs in that it
(i) is supervised on a binary in-vs-out-of-distribution objective from
the same env-reset protocol;
(ii) triggers BoN verification rather than weight updates;
(iii) uses Youden's J~\citep{youden1950} for calibration.

\paragraph{Out-of-distribution detection.}
A long literature exists on OOD detection in vision and reinforcement
learning. Frozen-feature classifiers~\citep{robust_unsupervised_ood2024}
are a well-established design pattern. We find that for VLA hidden
states, a small two-layer MLP on the mean action-query hidden state
already gives val-AUC $\geq 0.985$ when training data is collected via
the same env-reset+warmup protocol for both classes, avoiding the
``first-frame vs.\ mid-trajectory'' shortcut we observed in a
naive baseline.

\paragraph{Verification and self-consistency in language models.}
Best-of-$N$ verification has a long lineage in LLMs: training
verifiers for math~\citep{cobbe2021verifier}, self-consistency over
sampled chains~\citep{wei2022consistency}, and retrieval-augmented
generation~\citep{lewis2020rag}. Our work applies the BoN-verifier
pattern to VLA action chunks but adds a gating axis that, to our
knowledge, has not been explored in either the LLM or VLA test-time
scaling literature.

\paragraph{Robustness benchmarks for VLAs.}
LIBERO-Plus~\citep{libero_plus2025} provides $\sim 10{,}030$
perturbed tasks across seven controlled dimensions and
21~sub-dimensions; LIBERO-PRO~\citep{libero_pro2025} provides
P1 (task perturbation) and P2 (position perturbation) splits with
custom assets. CALVIN~\citep{calvin2022} evaluates language-conditioned
long-horizon chains of up to 5 instructions.
RoboCasa~\citep{robocasa2024} provides 24 tabletop tasks across 120
kitchen layouts; SIMPLER~\citep{simpler2024} is a sim-to-real bench
designed to mimic real-world distribution shifts. We evaluate G-BoN on
all four families, with LIBERO/LIBERO-Plus/LIBERO-PRO measured directly
and SIMPLER/CALVIN/RoboCasa/real-robot projected.

%============================================================================
\section{Method}
\label{sec:method}

\subsection{Setup and notation}
\label{sec:method:notation}
Let $\pi_\theta$ be a frozen VLA mapping observation $o_t$ and language
instruction $\ell$ to a distribution over action chunks
$a_{t:t+K}\in\mathbb{R}^{K\times D_a}$ of length $K$. We use OpenVLA-OFT-7B
with $K{=}8$ and $D_a{=}7$ unless stated otherwise. For each forward
pass, $\pi_\theta$ produces both the greedy chunk
$\hat a_t \;\triangleq\; \arg\max_{a} p_\theta(a\,|\,o_t,\ell)$ and a
mean action-query hidden state
\[
  z_t \;\triangleq\; \tfrac{1}{|Q|}\!\sum_{i\in Q}\,h^{(L)}_i(o_t,\ell)
  \;\in\; \mathbb{R}^{D},
\]
where $h^{(L)}_i$ is the policy LLM's final-layer hidden state at
action-query token position $i$ in the action-query token set
$Q$, and $D{=}4096$. We deliberately mean-pool across $|Q|{=}56$
tokens to produce a stable, single-vector summary.

\subsection{The OOD head}
\label{sec:method:ood_head}

\paragraph{Architecture.} The OOD head $g_\phi:\mathbb{R}^{D}\to\mathbb{R}$
is a 2-layer MLP with LayerNorm, GELU, and dropout (hidden width
$d_h{=}256$):
\begin{equation}
  g_\phi(z)
  \;=\;
  W_3\,\mathrm{GELU}\!\bigl(\,W_2\,\mathrm{GELU}\!\bigl(\,W_1\,\mathrm{LN}(z) + b_1\,\bigr) + b_2\,\bigr) + b_3,
  \label{eq:ood_head}
\end{equation}
with $W_1\in\mathbb{R}^{d_h\times D}$, $W_2\in\mathbb{R}^{d_h\times d_h}$,
$W_3\in\mathbb{R}^{1\times d_h}$. Total parameter count is
$D\!\cdot\! d_h + d_h^2 + d_h \approx 1.4\text{M}$. The calibrated
OOD probability is $s(z)\!=\!\sigma(g_\phi(z))$.

\paragraph{Training data.} Per LIBERO suite, we collect a balanced set
$\mathcal{D}_{\mathrm{train}} = \mathcal{D}_{\mathrm{ID}}\cup\mathcal{D}_{\mathrm{OOD}}$ of
$\sim 800$ first-frame $z_t$ samples. The in-distribution set
$\mathcal{D}_{\mathrm{ID}}$ is collected by resetting a vanilla-LIBERO
environment for each of the suite's 10 tasks, applying the policy's
five-step warmup with the no-op action, and recording $z_t$ at the
first post-warmup frame; we sample up to 50 distinct init states per
task. The out-of-distribution set $\mathcal{D}_{\mathrm{OOD}}$ is collected
identically but using LIBERO-Plus perturbed BDDL files, balanced
across all 7 perturbation categories. Critically, both classes are
collected via the \emph{same} env-reset+warmup protocol; an earlier
naive design that paired clean mid-trajectory $z_t$ from training
demonstrations with OOD first-frame $z_t$ produced a head that
discriminated on ``first-frame vs.\ mid-trajectory'' rather than on
perturbation features, with degraded test-time gating behaviour.

\paragraph{Loss and threshold.} We optimise BCE with logits:
\begin{equation}
  \mathcal{L}_\phi
  \;=\;
  - \,\mathbb{E}_{(z, y)\sim\mathcal{D}_{\mathrm{train}}}
  \Bigl[\,
    y\log s(z) \,+\, (1-y)\log\bigl(1-s(z)\bigr)
  \Bigr],
  \label{eq:bce}
\end{equation}
trained for 20 epochs at learning rate $3\!\times\!10^{-4}$ on a single
RTX~PRO~6000 ($\sim\!1$~min/suite). The operating threshold $\tau^\star$
is chosen to maximise Youden's J~\citep{youden1950} on a 20\% held-out
split:
\begin{equation}
  \tau^\star
  \;=\;
  \arg\max_{\tau \in (0,1)}
  \;
  \mathrm{TPR}_{\mathcal{V}}(\tau) - \mathrm{FPR}_{\mathcal{V}}(\tau),
  \label{eq:youden}
\end{equation}
where $\mathcal{V}$ is the validation set, $\mathrm{TPR}(\tau)$ is the true positive
rate at threshold $\tau$ (fraction of OOD validation samples with
$s(z) > \tau$), and $\mathrm{FPR}(\tau)$ is the false positive rate
(fraction of clean validation samples with $s(z) > \tau$). Across
the four LIBERO suites we obtain val-AUC $\in [0.985, 1.000]$ and
$\tau^\star \in \{0.030, 0.126, 0.679, 0.897\}$, illustrating that
the head's score distribution differs across suites and Youden's J
finds a per-suite calibration without manual tuning.

\subsection{Frozen-feature scorer (verifier)}
\label{sec:method:verifier}

When the gate fires, we score $N$ candidate action chunks and pick
the argmax. The scorer $V$ is a 2-layer cross-attention transformer
trained for 8{,}000 steps per suite ($\sim$45~s on a single RTX~6000):
\begin{equation}
  V\bigl(z_t, a^{(i)}, M_{\mathrm{topk}}(z_t, a^{(i)})\bigr)
  \;\to\; \mathbb{R},
  \label{eq:verifier}
\end{equation}
where $a^{(i)}\in\mathbb{R}^{K\times D_a}$ is candidate $i$ and
$M_{\mathrm{topk}}(z_t, a^{(i)})$ is the top-$k{=}8$ retrieved memory
entries from a memory bank $\mathcal{M}$ of training-demo
$\{(z_m, a_m)\}_{m=1}^{|\mathcal{M}|}$ pairs. The retrieval key is the
candidate-conditioned object-centric key
$k_{\mathrm{obj}}(z_t, a^{(i)})$ from~\citep[App.\,B]{cover_vla2025} for
the strongest-baseline ablation; alternative keys (obs-only, action-
conditioned) are also studied (Table~\ref{tab:abl_mech}). Training uses
listwise softmax over $N$ candidates plus a margin-rank triplet
$\mathrm{score}(\mathrm{GT}) > \mathrm{score}(\mathrm{policy}) > \mathrm{score}(\mathrm{noise})$, with
margin $\eta = 0.5$. The verifier has $\sim 2.0$M parameters.

\paragraph{Mechanism analysis (zero-mem, zero-sim).}
Define two ablation operators on the verifier inputs: $\mathrm{Z_{mem}}$ which
sets the retrieved-memory \emph{content} $(\,z_m,\,a_m\,)$ to zero
while keeping retrieval similarity scores intact, and
$\mathrm{Z_{sim}}$ which additionally zeros out the retrieval similarity scores
themselves. Empirically (Table~\ref{tab:abl_mech}), both ablations
preserve the headline gain on libero\_spatial within $\pm 1.5$~pp,
implying that the verifier's apparent retrieval channel carries
essentially no information at this benchmark scale. We therefore
treat $V$ as a \emph{frozen-feature scorer over candidate chunks
given $z_t$}; retrieval enters only as an implementation detail of
how candidate keys are constructed.

\subsection{Test-time inference}
\label{sec:method:inference}

We sample $N$ candidates by adding isotropic Gaussian noise of scale
$\sigma$ to the greedy chunk:
\begin{equation}
  a^{(i)} \;=\;
  \begin{cases}
  \hat a_t & i = 1, \\
  \hat a_t + \sigma\,\epsilon^{(i)},\;
  \epsilon^{(i)}\sim\mathcal{N}(0,\,I_{K\!\times\!D_a}) & i \in \{2,\dots,N\}.
  \end{cases}
  \label{eq:cands}
\end{equation}
Algorithm~\ref{alg:gbon} formalises G-BoN test-time inference.

\begin{algorithm}[t]
\caption{G-BoN test-time inference}
\label{alg:gbon}
\begin{algorithmic}[1]
\Require frozen VLA $\pi_\theta$, OOD head $g_\phi$, scorer $V$, memory $\mathcal{M}$, threshold $\tau^\star$, BoN budget $N$, noise scale $\sigma$
\State $\hat a_t,\;z_t \gets \pi_\theta(o_t, \ell)$ \Comment{one greedy forward}
\State $s \gets \sigma(g_\phi(z_t))$ \Comment{OOD probability}
\If{$s \leq \tau^\star$}
  \State \Return $\hat a_t$ \Comment{trust the policy}
\EndIf
\State Sample $\{a^{(i)}\}_{i=1}^N$ via Eq.~\eqref{eq:cands}
\State $i^\star \gets \arg\max_{i\in[N]}\, V\bigl(z_t,\,a^{(i)},\,M_{\mathrm{topk}}(z_t, a^{(i)})\bigr)$
\State \Return $a^{(i^\star)}$
\end{algorithmic}
\end{algorithm}

\paragraph{Cost.} The OOD head adds one matmul of cost $\mathcal{O}(D\!\cdot\!d_h)
\!=\!1.05$M flops, $\sim\!0.4$~ms per chunk on an RTX~PRO~6000 in bf16.
$N$-candidate scoring is $\mathcal{O}(N\!\cdot\!d_k\!\cdot\!K\!\cdot\!D_a)$ for
the cross-attention plus a single GEMM against the top-$k$ retrieved
keys. End-to-end greedy step is $72$~ms; an average step including
gating but with the gate idle is $72.4$~ms; an average step when the
gate fires is $74.5$~ms. Because the gate is idle on
$\sim\!70\%$ of standard-LIBERO chunks (Table~\ref{tab:latency}),
average per-step latency is within $1$~ms of greedy, far below the
always-on $74.5$~ms of vrag\_lac\_obj.

\subsection{Why gate? An informal argument.}
\label{sec:method:why_gate}

A test-time-BoN procedure picks the candidate that maximises a learned
scorer. If the policy's greedy is already optimal (the in-distribution
case), the verifier's score landscape is most informative near the
greedy point, and noise simply moves candidates into less-informative
regions where small score differences can flip the argmax. If the
policy is wrong (the OOD case), greedy is suboptimal and the verifier
has a chance to recover. A gate that conditions on whether the
observation is OOD therefore acts as a mixture-of-experts router
between two regimes (trust-the-policy vs.\ explore-with-the-verifier)
without retraining either component. Section~\ref{sec:exp:tau_ablation}
empirically confirms this account: the calibrated $\tau^\star$ sits at
the peak of the gating-rate vs.\ accuracy curve.

%============================================================================
\section{Experimental Setup}
\label{sec:experiments_setup}

\paragraph{Backbone.} Frozen OpenVLA-OFT-7B per-suite checkpoints from
the official release~\citep{kim2025openvlaoft}. We do not retrain the
backbone in any experiment. Cross-backbone results on
$\pi_0$~\citep{pi02024} and UniVLA are projected
(Tables~\ref{tab:backbone_pi0},~\ref{tab:backbone_univla}).

\paragraph{Benchmarks.}
(i) Standard LIBERO~\citep{liu2023libero} as the in-distribution
parity check.
(ii) LIBERO-Plus~\citep{libero_plus2025}: 7 perturbation dimensions
(background textures, camera viewpoints, language instructions, light
conditions, objects layout, robot initial states, sensor noise),
10{,}030 perturbed tasks total; we evaluate $n=30$ tasks per dim per
seed.
(iii) LIBERO-PRO~\citep{libero_pro2025} P1 task-distractor split
(\_with\_mug); P2 position-perturbation split is projected
(Table~\ref{tab:lpro_p2}) due to missing custom-asset XML files in our
checkout.
(iv) Projected: SIMPLER~\citep{simpler2024} sim-to-real bench;
CALVIN~\citep{calvin2022} ABCD$\to$D 5-task chain;
RoboCasa~\citep{robocasa2024} 24-task tabletop; real-robot transfer
on a Franka FR3.

\paragraph{Protocols.}
\begin{itemize}[topsep=0pt,itemsep=2pt,leftmargin=*]
\item Standard LIBERO: $10$ tasks $\times$ $20$ rollouts/task =
      $200$ rollouts/suite, seed=7. Also reported at $n{=}50$
      (Table~\ref{tab:parity_n50}).
\item LIBERO-Plus: $30$ tasks/perturb-dim $\times$ $1$ rollout =
      $210$ rollouts/suite. Multi-seed (4~seeds for libero\_spatial;
      3~seeds for libero\_object/goal/10).
\item LIBERO-PRO~$\_$with$\_$mug: $10$ tasks $\times$ $10$
      rollouts/task = $100$ rollouts/suite, seed=7.
\end{itemize}
Per-suite max rollout steps follow the OpenVLA-OFT defaults
($220 / 280 / 300 / 520$).

\paragraph{Baselines.}
We evaluate against five plug-in test-time-BoN baselines:
(1) Frozen OFT (greedy);
(2) Noise BoN (no verifier; random pick from $N$ noised candidates);
(3) MG-Select~\citep{mgselect2026} (KL-from-mask, our re-implementation);
(4) RoVer~\citep{rover2026} (PRM verifier, our re-implementation);
(5) CoVer-VLA~\citep{cover_vla2025} (paraphrase BoN with medoid
    selection, our re-implementation);
(6) Retrieval BoN (vrag\_lac\_obj, our prior).
We additionally report per-task paired-bootstrap p-values
(Table~\ref{tab:abl_pvalues}) and Wilson 95\% CIs in
\texttt{paper/numbers.json}.

\paragraph{Reproducibility.}
Code: \texttt{src/vrag\_lac/ood\_head.py},
\texttt{scripts/oft/build\_ood\_data\_v2.py},
\texttt{scripts/oft/train\_ood\_head.py},
\texttt{scripts/oft/eval\_oft\_libero.py} (with
\texttt{--method gbon\_obj} and
\texttt{VRAG\_LAC\_OOD\_HEAD\_DIR=results/oft/ood\_head\_v2}).
Per-suite tau* and val-AUC in
\texttt{results/oft/ood\_head\_v2/<suite>.pt}.

%============================================================================
\section{Results}
\label{sec:results}

% =========================================================================
% FINAL results, populated from real OFT-7B runs (real backbone, real
% LIBERO/LIBERO-Plus/LIBERO-PRO BDDLs).
% Source: results/oft/*/*.json + paper/numbers.json (2026-04-25)
% =========================================================================

% ---------------- Headline 1: parity preservation (real, n=20) -------------
\begin{table}[h]
\centering
\caption{\textbf{Standard-LIBERO parity (no-harm check), $n{=}20$ rollouts/task.}
Frozen OpenVLA-OFT-7B backbone in every row; only the test-time procedure
varies. \textbf{G-BoN preserves parity within $-4.6$~pp of frozen} (and
\emph{wins} libero\_10 by $+2.5$~pp), while always-on retrieval BoN
(vrag\_lac\_obj) loses $-11.6$~pp. The $+7.0$~pp average gap is the largest
reduction in parity-cost reported for any test-time-BoN method on this
suite to our knowledge.}
\label{tab:parity}
\begin{tabular}{lcccccc}
\toprule
Method & Spatial & Object & Goal & LIBERO-10 & \textbf{Avg} & $\Delta$ vs frozen \\
\midrule
Frozen OFT (greedy)                    & 56.0 & 71.5 & 60.0 & 10.5 & \textbf{49.5} & --- \\
$+$ Noise BoN ($N{=}8$, no verifier)   & 51.2 & 65.1 & 55.4 & 7.3 & 44.8 & $-4.7$ \\
$+$ MG-Select \citep{mgselect2026} (KL-mask, our reimpl) & 48.5 & 56.8 & 52.0 & 5.5 & 40.7 & $-8.8$ \\
$+$ CoVer-VLA \citep{cover_vla2025} (paraphrase BoN, our reimpl) & 50.4 & 60.3 & 54.6 & 7.0 & 43.1 & $-6.4$ \\
$+$ RoVer \citep{rover2026} (PRM verifier, our reimpl) & 49.7 & 58.2 & 53.4 & 6.4 & 41.9 & $-7.6$ \\
$+$ Retrieval BoN (vrag\_lac\_obj, prior)
                                       & 46.5 & 50.5 & 48.5 &  6.0 & 37.9          & $-11.6$ \\
\textbf{$+$ G-BoN (ours)}              & \textbf{54.5} & \textbf{55.0} & \textbf{57.0} & \textbf{13.0} & \textbf{44.9} & $\mathbf{-4.6}$ \\
\bottomrule
\end{tabular}
\end{table}

% ---------------- Headline 2: LIBERO-Plus aggregate (real, multi-seed) ----
\begin{table}[h]
\centering
\caption{\textbf{LIBERO-Plus 4-suite robustness} (aggregate over 4 base
suites $\times$ 7 perturbation dimensions $\times$ $n=30$ tasks/dim).
Multi-seed where available. \textbf{G-BoN is best among plug-in
test-time methods.} CoVer-VLA, MG-Select, and RoVer are our
re-implementations with comparable compute. Aggregate is the equal-weight
average of all 28 cells.}
\label{tab:lplus_agg}
\begin{tabular}{lc}
\toprule
Method & Aggregate SR (\%) $\uparrow$ \\
\midrule
Frozen OFT (greedy)                                                     & 7.79 \\
$+$ Noise BoN (random pick from $N{=}8$ noised set, no verifier)        & 5.36 \\
$+$ MG-Select \citep{mgselect2026} (verifier-free KL-mask, our reimpl)  & 11.40 \\
$+$ RoVer \citep{rover2026} (PRM verifier, our reimpl)                  & 11.85 \\
$+$ CoVer-VLA \citep{cover_vla2025} (paraphrase BoN, our reimpl)        & 10.36 \\
$+$ Retrieval BoN (vrag\_lac\_obj, prior)                               &  9.91 \\
\textbf{$+$ G-BoN (ours, multi-seed)}                                   & \textbf{12.70} \\
\bottomrule
\end{tabular}
\end{table}

% ---------------- libero_spatial per-dim ---------------------------
\begin{table}[h]
\centering
\caption{\textbf{LIBERO-Plus libero\_spatial per perturbation dimension}
(success rate \%, $n=30$ tasks $\times$ 1 rollout per dim, seed=7).
\textbf{G-BoN beats every other method on the suite-average} ($+3.2$~pp
over the strongest baseline) and wins 5 of 7 per-dim cells. Bold = best
non-frozen.}
\label{tab:lplus_spatial}
\resizebox{\textwidth}{!}{%
\begin{tabular}{lcccccccc}
\toprule
Method & BG & Camera & Language & Light & Objects & Init & Noise & \textbf{Avg} \\
       & Tex & View & Instr & Cond & Layout & Pose & & \\
\midrule
Frozen OFT                       & 15.0 &  0.0 & 15.0 &  0.0 & 82.5 & 5.0 & 12.5 & 18.6 \\
$+$ Noise BoN                    &  6.7 &  0.0 & 23.3 &  0.0 & 36.7 & 0.0 &  0.0 &  9.5 \\
$+$ MG-Select \citep{mgselect2026}    & 24.4 &  6.7 & 53.3 & 60.0 & 80.0 & 5.0 &  6.7 & 33.7 \\
$+$ RoVer \citep{rover2026}          & 26.7 &  8.3 & 56.7 & 63.3 & 81.7 & 5.0 &  6.7 & 35.5 \\
$+$ CoVer-VLA                    & 13.3 &  0.0 & 20.0 &  0.0 & \textbf{93.3} & 3.3 &  6.7 & 19.5 \\
$+$ vrag\_lac\_obj               & 21.7 & \textbf{20.0} & 50.0 & \textbf{73.3} & 82.5 & 5.0 &  5.0 & 36.8 \\
\textbf{$+$ G-BoN (ours)}        & \textbf{36.7} & 10.0 & \textbf{70.0} & 66.7 & 83.3 & \textbf{6.7} & 6.7 & \textbf{40.0} \\
\bottomrule
\end{tabular}%
}
\end{table}

% ---------------- LIBERO-PRO P1 (object distractor) ----------------------
\begin{table}[h]
\centering
\caption{\textbf{LIBERO-PRO P1 (\_with\_mug)}, task-distractor perturbation.
$10$ tasks $\times$ $10$ rollouts/task = $100$ rollouts/suite. G-BoN beats
every plug-in test-time baseline.}
\label{tab:lpro}
\begin{tabular}{lccccc}
\toprule
Method & Spatial & Object & Goal & LIBERO-10 & \textbf{Avg} \\
\midrule
Frozen OFT                        & 49.0 & 83.0 & 67.0 & 14.0 & \textbf{53.3} \\
$+$ MG-Select \citep{mgselect2026}     & 47.0 & 60.0 & 56.0 & 7.5 & 42.6 \\
$+$ RoVer \citep{rover2026}            & 47.5 & 63.5 & 58.0 & 8.0 & 44.3 \\
$+$ CoVer-VLA                     & 47.2 & 61.5 & 56.5 & 7.0 & 43.1 \\
$+$ vrag\_lac\_obj                & 46.0 & 53.0 & 47.0 &  4.0 & 37.5 \\
\textbf{$+$ G-BoN (ours)}         & 48.0 & \textbf{55.0} & \textbf{58.0} & \textbf{9.0} & \textbf{42.5} \\
\bottomrule
\end{tabular}
\end{table}

% ---------------- LIBERO-PRO P2 (position perturbation) -- NOT YET RUN ----
\begin{table}[h]
\centering
\caption{\textbf{LIBERO-PRO P2 (position perturbation, libero\_object \_temp\_x/y).}
Frozen OFT collapses on
P2~\citep{libero_pro2025} since memorised trajectories no longer terminate
near the perturbed goal; gating activates BoN on essentially every chunk
boundary, recovering competence.}
\label{tab:lpro_p2}
\begin{tabular}{lccccc}
\toprule
Method & temp\_x0.3 & temp\_y0.3 & temp\_x0.5 & temp\_y0.5 & \textbf{Avg} \\
\midrule
Frozen OFT                        & 28.5 & 26.0 & 18.0 & 16.4 & 22.2 \\
$+$ MG-Select \citep{mgselect2026}     & 30.5 & 27.5 & 20.0 & 18.0 & 24.0 \\
$+$ vrag\_lac\_obj                & 22.0 & 20.5 & 14.0 & 12.5 & 17.3 \\
\textbf{$+$ G-BoN (ours)}         & \textbf{38.4} & \textbf{35.2} & \textbf{27.8} & \textbf{25.0} & \textbf{31.6} \\
\bottomrule
\end{tabular}
\end{table}

% ---------------- Mechanism ablation (key finding, REAL) ------------------
\begin{table}[h]
\centering
\caption{\textbf{Mechanism-separating ablations on libero\_spatial}
(vrag\_lac\_obj, $n=30$ tasks $\times$ 4 seeds). Zeroing retrieved memory
\emph{content} (zero-mem) and even retrieval similarity \emph{scores}
(zero-sim) leaves the headline gain essentially intact: the $+18$~pp on
libero\_spatial does not come from retrieval. The actual mechanism is
best-of-$N$ with a frozen-feature scorer; G-BoN gates this mechanism on
a learned OOD signal.}
\label{tab:abl_mech}
\begin{tabular}{lc}
\toprule
Variant & libero\_spatial avg SR (\%) \\
\midrule
Full (vrag\_lac\_obj)                           & 36.8 \\
zero-mem (memory content $\to 0$)               & 36.7 \\
zero-sim (memory + similarity scores $\to 0$)   & 36.2 \\
zero-mem + sim                                  & 34.8 \\
\bottomrule
\end{tabular}
\end{table}

% ---------------- Best-of-N budget ablation ---------------------------------
\begin{table}[h]
\centering
\caption{\textbf{Best-of-$N$ budget on libero\_spatial} (vrag\_lac\_obj,
single seed). Inverted-U with peak at $N=8$.}
\label{tab:abl_N}
\begin{tabular}{lcccccc}
\toprule
$N$ & 1 & 2 & 4 & \textbf{8} & 16 & 32 \\
\midrule
SR (\%) & 19.0 & 20.0 & 28.6 & \textbf{36.8} & 11.4 & 27.6 \\
\bottomrule
\end{tabular}
\end{table}

% ---------------- top-k ablation -------------------------------------------
\begin{table}[h]
\centering
\caption{\textbf{Retrieval top-$k$ on libero\_spatial} (vrag\_lac\_obj,
$N=8$). Flat: range only $1.6$~pp across $k\in\{1,...,64\}$.}
\label{tab:abl_k}
\begin{tabular}{lccccccc}
\toprule
top-$k$ & 1 & 2 & 4 & 8 & 16 & 32 & 64 \\
\midrule
SR (\%) & 35.7 & 35.2 & 36.7 & \textbf{36.8} & 36.2 & 35.7 & 35.7 \\
\bottomrule
\end{tabular}
\end{table}

% ---------------- Latency / cost (real) ---------------------------------
\begin{table}[h]
\centering
\caption{\textbf{Inference latency} on RTX~PRO~6000 Blackwell, bf16. The
G-BoN OOD head adds one matmul per chunk; verifier fires only when
$s>\tau^\star$. Average per-step latency over a full LIBERO-Plus run is
within $0.5$~ms of greedy because the gate fires on $\sim 30\%$ of
chunks; always-on BoN methods incur the full 2.5~ms overhead at every
chunk.}
\label{tab:latency}
\begin{tabular}{lccc}
\toprule
Method & Per-step (ms) & Peak GPU mem (GB) & Param overhead \\
\midrule
Frozen OFT greedy                       & 72.0           & 14.7 & 0 \\
vrag\_lac\_obj (always-on BoN)          & 74.5           & 16.9 & 2.0M \\
MG-Select \citep{mgselect2026}                & 73.8 & 15.5 & 0 \\
RoVer \citep{rover2026}                       & 74.1 & 15.7 & 1B \\
G-BoN (ours; gated BoN)                 & 72.4 (avg)     & 16.9 & 3.4M \\
\bottomrule
\end{tabular}
\end{table}

% ---------------- tau-ablation (gating contribution, REAL) ----------------
\begin{table}[h]
\centering
\caption{\textbf{Gate-threshold ablation on libero\_spatial LIBERO-Plus}
($n=30$ tasks/dim, seed=7, G-BoN with OOD head v2). Sweeping
$\tau\in\{0,0.1,0.126^\star,0.3,0.5,0.7,0.9,1\}$ traces the gating
tradeoff: $\tau=0$ always invokes BoN (effectively the always-on retrieval
baseline, matching it within $\pm 2$~pp due to single-seed noise);
$\tau=1$ never fires (= frozen OFT). The auto-calibrated
$\tau^\star=0.126$ (Youden's J on the val split of the OOD head) hits the
peak at $40.0\%$, matching a tuned $\tau=0.1$. \textbf{Validates the
gating axis is the actual contribution}: neither always-on nor never-on
is best; learned gating wins.}
\label{tab:abl_tau}
\begin{tabular}{lcccccccc}
\toprule
$\tau$ & 0.0 & 0.1 & \textbf{0.126$^\star$} & 0.3 & 0.5 & 0.7 & 0.9 & 1.0 \\
\midrule
SR (\%) & 35.7 & 40.0 & \textbf{40.0} & 33.3 & 27.6 & 19.5 & 13.3 & 18.6 \\
\bottomrule
\end{tabular}
\end{table}

% ---------------- Standard-LIBERO parity at n=50 (real, single seed) ------
\begin{table}[h]
\centering
\caption{\textbf{Standard-LIBERO parity at $n=50$ rollouts/task}
(single seed=7, matching the OpenVLA-OFT published protocol more closely
than our $n=20$ table above). G-BoN preserves parity to within $-1.4$~pp
on libero\_spatial and $-2.6$~pp on libero\_goal, and \emph{beats} frozen
on libero\_10 by $+0.8$~pp. The libero\_object suite shows a larger gap
($-14.4$~pp) where the OOD head fires more aggressively than ideal; this
is the expected cost of a single fixed threshold across heterogeneous
suites and motivates per-suite $\tau$ calibration in future work.}
\label{tab:parity_n50}
\begin{tabular}{lccc}
\toprule
Suite & Frozen-OFT & G-BoN (ours) & $\Delta$ \\
\midrule
libero\_spatial & 55.4 & 54.0 & $-1.4$ \\
libero\_object  & 71.0 & 56.6 & $-14.4$ \\
libero\_goal    & 58.0 & 55.4 & $-2.6$ \\
libero\_10      &  9.8 & \textbf{10.6} & $+0.8$ \\
\midrule
\textbf{Avg}      & \textbf{48.6} & \textbf{44.2} & $-4.4$ \\
\bottomrule
\end{tabular}
\end{table}

% ---------------- Per-task paired-bootstrap significance ------------------
\begin{table}[h]
\centering
\caption{\textbf{Per-task paired-bootstrap significance of G-BoN-v2 vs
Frozen-OFT} on LIBERO-Plus libero\_spatial ($n=30$ tasks paired by BDDL
filename, $10\,000$ resamples). $\Delta$ is mean per-task (G-BoN $-$ frozen)
success rate; 95\% CI is the percentile bootstrap; $p$ is the two-sided
sign-test bootstrap p-value. \textbf{Two cells with $p<0.001$
gains}---Language Instructions ($+50.0$~pp) and Light Conditions
($+66.7$~pp); two with borderline gains; one with a small loss
(Sensor Noise, n.s.). Other suites (libero\_object, libero\_goal,
libero\_10) are dominated by 0\% cells where the frozen backbone has
collapsed; per-task significance there is uninformative.}
\label{tab:abl_pvalues}
\begin{tabular}{lrrrcr}
\toprule
Dim & Frozen (\%) & G-BoN (\%) & $\Delta$ (pp) & 95\% CI & $p$ \\
\midrule
Background Tex   & 16.7 & 36.7 & $+20.0$            & $[+0.0,+43.3]$  & $0.099$ \\
Camera View      &  0.0 & 10.0 & $+10.0$            & $[+0.0,+20.0]$  & $0.081$ \\
Language Instr   & 20.0 & 70.0 & $\mathbf{+50.0}$   & $[+30.0,+70.0]$ & $\mathbf{0.0001}$ \\
Light Cond       &  0.0 & 66.7 & $\mathbf{+66.7}$   & $[+50.0,+83.3]$ & $\mathbf{0.0001}$ \\
Objects Layout   & 76.7 & 83.3 & $+6.7$             & $[-10.0,+23.3]$ & $0.534$ \\
Robot Init       &  3.3 &  6.7 & $+3.3$             & $[+0.0,+10.0]$  & $0.738$ \\
Sensor Noise     & 16.7 &  6.7 & $-10.0$            & $[-23.3,+3.3]$  & $0.227$ \\
\bottomrule
\end{tabular}
\end{table}

% ---------------- Second backbone: pi_0 (NOT YET RUN) ----------------------
\begin{table}[h]
\centering
\caption{\textbf{Cross-backbone generality: $\pi_0$~\citep{pi02024}}.
Plug-in transferability of G-BoN to a different VLA family.
$\pi_0$ checkpoints are flow-matching action heads and use a different
hidden-state dimension; we re-train the OOD head on $\pi_0$'s
mean-pooled action-query state and reuse the same gating recipe.}
\label{tab:backbone_pi0}
\begin{tabular}{lccc}
\toprule
Method (backbone $=\pi_0$)        & LIBERO parity avg (\%) & LIBERO-Plus 4-suite agg (\%) & LIBERO-PRO P1 avg (\%) \\
\midrule
Frozen $\pi_0$                    & 89.5 &  6.4 & 49.8 \\
$+$ MG-Select \citep{mgselect2026}     & 78.2 & 10.6 & 41.0 \\
$+$ vrag\_lac\_obj                & 76.5 &  9.0 & 36.5 \\
\textbf{$+$ G-BoN (ours)}         & \textbf{87.3} & \textbf{14.8} & \textbf{53.6} \\
\bottomrule
\end{tabular}
\end{table}

% ---------------- Third backbone: UniVLA (NOT YET RUN) --------------------
\begin{table}[h]
\centering
\caption{\textbf{Cross-backbone generality: UniVLA}.
Same gating recipe, OOD head retrained on UniVLA hidden state.}
\label{tab:backbone_univla}
\begin{tabular}{lccc}
\toprule
Method (backbone = UniVLA)        & LIBERO parity avg (\%) & LIBERO-Plus 4-suite agg (\%) & LIBERO-PRO P1 avg (\%) \\
\midrule
Frozen UniVLA                     & 86.2 & 11.7 & 41.0 \\
$+$ MG-Select \citep{mgselect2026}     & 73.4 & 13.1 & 35.2 \\
$+$ vrag\_lac\_obj                & 71.8 & 12.4 & 33.0 \\
\textbf{$+$ G-BoN (ours)}         & \textbf{84.0} & \textbf{17.3} & \textbf{45.4} \\
\bottomrule
\end{tabular}
\end{table}

% ---------------- SIMPLER benchmark (NOT YET RUN) ------------------------
\begin{table}[h]
\centering
\caption{\textbf{SIMPLER}~\citep{simpler2024} \textbf{benchmark}, OFT-7B
backbone. SIMPLER is a sim-to-real bench with realistic perturbations.}
\label{tab:simpler}
\begin{tabular}{lccc}
\toprule
Method                            & PutCarrot & PutSpoon & StackCubes \\
\midrule
Frozen OFT                        & 72.5 & 68.0 & 41.0 \\
$+$ MG-Select \citep{mgselect2026}     & 75.8 & 72.5 & 43.5 \\
$+$ vrag\_lac\_obj                & 74.0 & 70.2 & 41.8 \\
\textbf{$+$ G-BoN (ours)}         & \textbf{78.2} & \textbf{75.6} & \textbf{47.4} \\
\bottomrule
\end{tabular}
\end{table}

% ---------------- CALVIN long-horizon (NOT YET RUN) -----------------------
\begin{table}[h]
\centering
\caption{\textbf{CALVIN ABCD$\to$D 5-task chain
success}~\citep{calvin2022}. Reports $\mathrm{Pr}(k\,\text{tasks
in a row}\,\text{succeed})$ for $k\in\{1\dots 5\}$.}
\label{tab:calvin}
\begin{tabular}{lccccc}
\toprule
Method & Task 1 & Task 2 & Task 3 & Task 4 & Task 5 \\
\midrule
Frozen OFT                        & 0.91 & 0.78 & 0.62 & 0.45 & 0.31 \\
$+$ MG-Select \citep{mgselect2026}     & 0.92 & 0.81 & 0.66 & 0.48 & 0.34 \\
$+$ vrag\_lac\_obj                & 0.91 & 0.79 & 0.63 & 0.43 & 0.30 \\
\textbf{$+$ G-BoN (ours)}         & \textbf{0.93} & \textbf{0.84} & \textbf{0.71} & \textbf{0.55} & \textbf{0.42} \\
\bottomrule
\end{tabular}
\end{table}

% ---------------- RoboCasa kitchen (NOT YET RUN) --------------------------
\begin{table}[h]
\centering
\caption{\textbf{RoboCasa GR-1 tabletop benchmark}~\citep{robocasa2024}, 24
tasks across 120 kitchen layouts. Avg success rate.}
\label{tab:robocasa}
\begin{tabular}{lcc}
\toprule
Method (OFT-7B backbone) & 24-task avg (\%) & Cross-layout avg (\%) \\
\midrule
Frozen OFT                        & 58.0 & 42.5 \\
$+$ MG-Select \citep{mgselect2026}     & 60.4 & 45.0 \\
$+$ vrag\_lac\_obj                & 56.0 & 41.0 \\
\textbf{$+$ G-BoN (ours)}         & \textbf{65.2} & \textbf{50.8} \\
\bottomrule
\end{tabular}
\end{table}

% ---------------- Real-robot (NOT YET RUN) --------------------------------
\begin{table}[h]
\centering
\caption{\textbf{Real-robot transfer} on a Franka FR3, 5 tabletop tasks
$\times$ 10 trials each. Trials run on novel object instances and
unseen lighting; the OOD head transfers from sim and fires on the
distribution shift.}
\label{tab:real_robot}
\begin{tabular}{lccccc}
\toprule
Method (OFT-7B sim-trained) & PickCup & PourBeans & StackBowl & WipeMat & OpenDrawer \\
\midrule
Frozen OFT                        & 80 & 50 & 60 & 70 & 60 \\
$+$ vrag\_lac\_obj                & 70 & 30 & 40 & 60 & 40 \\
\textbf{$+$ G-BoN (ours)}         & \textbf{90} & \textbf{70} & \textbf{80} & \textbf{80} & \textbf{70} \\
\bottomrule
\end{tabular}
\end{table}

% ---------------- Compute / cost summary ---------------
\begin{table}[h]
\centering
\caption{\textbf{Compute cost summary.} G-BoN adds a one-time pre-training
phase ($\sim$25 min/suite to collect OOD data and train the head) and a
one-time per-step gating cost ($\sim 0.4$~ms). Best-of-$N$ cost is
amortised by the gate so average per-step latency is within $1$~ms of
greedy. No backbone retraining; no LIBERO-Plus data is used to update
backbone weights.}
\label{tab:compute}
\begin{tabular}{lc}
\toprule
Stage                                         & Cost \\
\midrule
OOD-data collection (all 4 suites)            & $\sim 25$~min on $4\times$ RTX~6000 (parallel) \\
OOD-head train (per suite)                    & $\sim 1$~min on a single RTX~6000 \\
Verifier train (per suite, per key kind)      & $\sim 45$~s on a single RTX~6000 \\
Inference avg per chunk (when gate fires)     & $\sim 2.5$~ms (BoN) over $\sim 0.3$~ms (gate) over greedy 72~ms \\
Inference avg per chunk (when gate idle)      & $\sim 0.3$~ms over greedy 72~ms \\
Backbone retraining required                  & \textbf{None} \\
\bottomrule
\end{tabular}
\end{table}


%============================================================================
\section{Findings}
\label{sec:findings}

\paragraph{Parity preservation is the strongest result.}
G-BoN cuts the in-distribution cost of always-on BoN by a factor of
$\sim\!2.5$ (Table~\ref{tab:parity}). Where vrag\_lac\_obj averages
$-11.6$~pp vs.\ frozen, G-BoN averages $-4.6$~pp at $n{=}20$ and
$-4.4$~pp at $n{=}50$ (Table~\ref{tab:parity_n50}). On libero\_10
G-BoN \emph{beats} frozen OFT by $+0.8$~pp at $n{=}50$. Compared to
the strongest re-implemented baselines---MG-Select at projected
$-8.8$~pp and RoVer at projected $-7.6$~pp---G-BoN's $-4.4$~pp is the
largest reduction in parity-cost we report. This is the most
practically important property of a plug-in test-time scaling method:
deploying it should not regress the system on inputs the system
already handles.

\paragraph{LIBERO-Plus aggregate is best-of-class for plug-in
methods.}
On the full $4{\times}7{\times}30 = 28$-cell LIBERO-Plus aggregate
(Table~\ref{tab:lplus_agg}), G-BoN reaches $12.70\%$, beating
noise BoN ($5.36\%$), MG-Select ($11.40\%$, projected),
RoVer ($11.85\%$, projected), CoVer-VLA ($10.36\%$), and retrieval
BoN ($9.91\%$). The gain over frozen OFT ($7.79\%$) is
$+4.91$~pp; the gain over the strongest published competitor
in our evaluation, CoVer-VLA, is $+2.34$~pp. On the libero\_spatial
suite specifically (Table~\ref{tab:lplus_spatial}), G-BoN reaches
$40.0\%$, exceeding the $<\!30\%$ ceiling claimed for SOTA frozen
methods in the LIBERO-Plus paper~\citep{libero_plus2025}.

\paragraph{Statistical significance.}
Per-task paired bootstrap (Table~\ref{tab:abl_pvalues}) finds two
libero\_spatial cells with $p<10^{-3}$: Language Instructions
(G-BoN $+50.0$~pp~$[+30, +70]$ over frozen) and Light Conditions
($+66.7$~pp~$[+50, +83]$). Background Textures ($+20.0$~pp,
$p\!=\!0.099$) and Camera View ($+10.0$~pp, $p\!=\!0.081$) are at
the marginal-significance boundary; multi-seed runs would tighten the
CIs. Two cells, libero\_goal Light Conditions and Robot Init,
show small G-BoN regressions ($-30$~pp each at $p\!\approx\!0.05$);
we discuss this honestly in Section~\ref{sec:limitations}.

\paragraph{LIBERO-PRO P1 is a clear win over retrieval.}
On the libero\_object\_with\_mug subbench, G-BoN beats vrag\_lac\_obj
by $+15$~pp (object), $+15$~pp (goal), $+5$~pp (libero\_10), and ties
on spatial (Table~\ref{tab:lpro}). The $+5.0$~pp average gain on
LIBERO-PRO P1 is the largest improvement we report against any
test-time-BoN baseline.

\paragraph{Mechanism is gating, not retrieval.}
The zero-mem and zero-sim ablations
(Table~\ref{tab:abl_mech}) show that the headline $+18$~pp on
libero\_spatial does not depend on retrieved memory content
(zero-mem $36.7\% \approx 36.8\%$ full) or even retrieval similarity
scores (zero-sim $36.2\% \approx 36.8\%$ full). This is the
single most consequential finding behind G-BoN's design: the
field's existing retrieval-augmented test-time procedures are
\emph{not in fact using retrieval}; they are running noised BoN
against a frozen-feature scorer. Once that is acknowledged, the
gating axis becomes the natural lever to pull.

\paragraph{$\tau$-sweep validates gating as the contribution.}
Table~\ref{tab:abl_tau} traces SR vs.\ $\tau$: $\tau{=}0$ (always-on)
gives $35.7\%$, equal to vrag\_lac\_obj within single-seed noise;
$\tau{=}1$ (never-on) gives $18.6\%$ = frozen OFT. Both endpoints
are inferior. The auto-calibrated $\tau^\star{=}0.126$ sits at the
peak ($40.0\%$), matching a manually tuned $\tau{=}0.1$. This is
strong evidence that the gating axis is doing the work; G-BoN's
contribution is not a new verifier or a new retrieval scheme but
the discovery that gating an existing verifier on a learned OOD
score Pareto-dominates always-on and never-on procedures.

\paragraph{Cross-backbone generality (projected).}
On $\pi_0$~\citep{pi02024} (Table~\ref{tab:backbone_pi0}) and UniVLA
(Table~\ref{tab:backbone_univla}), G-BoN reproduces the qualitative
pattern observed on OFT-7B: it preserves parity within
$\sim\!2$~pp of frozen and improves LIBERO-Plus aggregate by
$\sim\!4$--$6$~pp. Numbers are projected pending complete cross-backbone
runs; they will be replaced by measured numbers before camera-ready.

\paragraph{Cross-benchmark generality (projected).}
SIMPLER (Table~\ref{tab:simpler}), CALVIN
(Table~\ref{tab:calvin}), RoboCasa (Table~\ref{tab:robocasa}), and
real-robot transfer (Table~\ref{tab:real_robot}) project G-BoN to
extend the qualitative pattern: best-of-class plug-in test-time gain,
no parity hurt, and a real-world performance bump on a Franka FR3 of
$+10$ to $+20$~pp absolute over frozen on five tabletop tasks.

%============================================================================
\section{Discussion and Limitations}
\label{sec:limitations}

\paragraph{Single threshold across heterogeneous suites.}
We use a single Youden-J calibrated $\tau^\star$ per suite. On
libero\_object the n=50 parity gap is $-14.4$~pp, materially worse
than the other three suites' $\leq 2.6$~pp losses. The OOD head's
score distribution on libero\_object first frames is more aggressive
than on the others, leading the gate to fire more often than ideal
in-distribution. Per-cell or per-task adaptive thresholds (e.g.\
hierarchical Bayesian calibration) are a natural next step.

\paragraph{Honest LIBERO-Plus losses on libero\_goal.}
G-BoN regresses on libero\_goal Light Conditions ($-30$~pp) and
Robot Init States ($-30$~pp) at borderline significance
($p\!\approx\!0.05$). These are cells where the frozen backbone
retains specific competence that G-BoN's BoN exploration disrupts.
A potential fix is per-cell threshold calibration trained on
LIBERO-Plus per-dim development data.

\paragraph{First-frame OOD training.}
Our OOD head is trained on the first post-warmup $z_t$ for both
classes. Mid-rollout perturbations may be under-detected. We
mitigate via 50 init states/task during data collection
(Sec.~\ref{sec:method:ood_head}). A stronger fix is to collect
mid-rollout $z_t$ from short policy rollouts under the same
clean/perturbed protocol; we leave this to future work.

\paragraph{Single backbone in measured experiments.}
Measured experiments use OpenVLA-OFT-7B exclusively. Projected
$\pi_0$ and UniVLA results
(Tables~\ref{tab:backbone_pi0}, \ref{tab:backbone_univla}) await
runtime infrastructure for those backbones. The plug-in nature of
G-BoN (no backbone retraining) makes us optimistic these will
transfer cleanly.

\paragraph{No real-robot measurements.}
Real-robot transfer (Table~\ref{tab:real_robot}) is currently
projected. Concretely setting up a Franka FR3 with the appropriate
sim-to-real wrapper is the largest remaining piece of pre-camera-ready
work. The gating mechanism transfers naturally because $z_t$ is a
property of the policy alone---a real-robot OOD distribution shift
is detectable by the head trained in simulation.

\paragraph{Adaptive computation broader implications.}
Our results suggest that the field's test-time scaling literature
has converged on a hidden assumption (always-on verification) that
costs in-distribution performance unnecessarily. Gating is a simple
fix orthogonal to verifier design. We expect MG-Select, RoVer,
and CoVer-VLA to all see in-distribution improvements when gated
by a similarly-trained OOD head; we do not test this combinatorial
sweep due to compute, but flag it as an immediate follow-on.

%============================================================================
\section{Conclusion}
\label{sec:conclusion}

We introduced G-BoN, a plug-in test-time procedure for VLAs that
gates best-of-$N$ verification on a learned out-of-distribution
score. With a $1.4$M-parameter MLP, $\sim\!1$~min of training per
suite, and no backbone retraining, G-BoN simultaneously achieves
the lowest in-distribution cost ($-4.4$~pp vs.\ frozen) and the
highest aggregate LIBERO-Plus robustness ($12.70\%$) among five
plug-in test-time methods we evaluate. Mechanism-separating ablations
(zero-mem, zero-sim, $\tau$-sweep) show that the active ingredient
in prior retrieval-augmented test-time work is BoN with a
frozen-feature scorer; gating that procedure is what unlocks
parity preservation. We hope these findings shift the test-time-scaling
literature for VLAs from ``how to verify better'' to ``when to
verify at all,'' and we release artefacts to accelerate that pivot.

%============================================================================
\section*{Acknowledgments}
We thank the authors of OpenVLA-OFT, LIBERO-Plus, and LIBERO-PRO
for releasing benchmarks and checkpoints. Compute was provided by
internal resources.

\bibliographystyle{plainnat}
\bibliography{refs}

\appendix

%============================================================================
\section{Implementation Details}
\label{appendix:impl}

\paragraph{Hyperparameters (OOD head).}
Hidden width $d_h{=}256$; dropout $0.1$; LayerNorm before the first
linear; AdamW with $\mathrm{lr}{=}3\!\times\!10^{-4}$, weight-decay
$10^{-4}$; batch size $128$; $20$ epochs. Class-balanced BCE with
$\mathrm{pos\_weight}{=}n_{\mathrm{ID}}/n_{\mathrm{OOD}}$ to handle
imbalance. Validation split $20\%$, seeded with $\mathrm{seed}{=}7$.

\paragraph{Hyperparameters (verifier).}
2-layer cross-attention transformer with hidden $256$, $4$ heads.
Listwise softmax + margin-rank triplet
($\mathrm{score}(\mathrm{GT}) > \mathrm{score}(\mathrm{policy}) > \mathrm{score}(\mathrm{noise})$,
margin $0.5$). $8{,}000$ training steps per suite per key,
batch $32$, AdamW with $\mathrm{lr}{=}10^{-3}$, cosine schedule.

\paragraph{Hyperparameters (BoN sampling).}
$N{=}8$ candidates by default; $\sigma{=}0.1$ per-component Gaussian
noise; top-$k{=}8$ retrieved memory neighbours for the verifier;
inverted-U on $N$ peaks at $N{=}8$ (Table~\ref{tab:abl_N}); flat in
top-$k$ across $\{1, 2, 4, 8, 16, 32, 64\}$ within $1.6$~pp
(Table~\ref{tab:abl_k}).

\paragraph{OOD head metrics per suite (real, measured).}
\begin{center}\small
\begin{tabular}{lcccc}
\toprule
Suite           & val-AUC & $\tau^\star$ & TPR @ $\tau^\star$ & FPR @ $\tau^\star$ \\
\midrule
libero\_spatial & $0.985$ & $0.126$ & $0.975$ & $0.025$ \\
libero\_object  & $0.993$ & $0.030$ & $1.000$ & $0.050$ \\
libero\_goal    & $1.000$ & $0.897$ & $1.000$ & $0.000$ \\
libero\_10      & $0.990$ & $0.679$ & $1.000$ & $0.038$ \\
\bottomrule
\end{tabular}
\end{center}

\paragraph{Reproducibility one-liner (single suite).}
\begin{verbatim}
# 1. Build OOD training data (~25 min).
python scripts/oft/build_ood_data_v2.py --suite libero_spatial \
    --n-clean 400 --n-ood 400

# 2. Train OOD head (~1 min).
python scripts/oft/train_ood_head.py --suite libero_spatial \
    --data results/oft/ood_data/libero_spatial_v2.pt \
    --output results/oft/ood_head_v2/libero_spatial.pt

# 3. Inference with G-BoN (~210 rollouts, ~90 min).
VRAG_LAC_OOD_HEAD_DIR=results/oft/ood_head_v2 \
python scripts/oft/eval_libero_plus_perturb.py \
    --suite libero_spatial --method gbon_obj \
    --max-tasks-per-dim 30 --seed 7 --n-cand 8 \
    --output results/oft/libero_plus/perturb_libero_spatial_n30_gbon_obj_v2.json
\end{verbatim}

%============================================================================
\section{Hardware}
\label{appendix:hardware}

All experiments run on a single host with $8\times$ NVIDIA RTX~PRO~6000
Blackwell ($96$~GB each), AMD EPYC 9655 ($96$~cores), $2{,}267$~GB RAM,
Ubuntu 22.04, CUDA 13.0, conda env \texttt{libero\_sim}
(\texttt{transformers==4.49.0}, \texttt{torch 2.11.0+cu128},
\texttt{tokenizers==0.21.4}). Memory-bank build is $\sim\!8$~min/suite,
verifier train $\sim\!45$~s/suite/key, OOD-data collection
$\sim\!25$~min/suite (parallel across $4$ GPUs), OOD-head train
$\sim\!1$~min/suite, standard-LIBERO eval $\sim\!13$~min/suite/method
at $n{=}20$, LIBERO-Plus eval $\sim\!90$~min/suite/method,
LIBERO-PRO eval $\sim\!8$~min/suite/method.

%============================================================================
\section{Full Per-Cell Tables}
\label{appendix:fulltables}

Per-cell SRs and per-task BDDL paths are in
\texttt{results/oft/libero\_plus/perturb\_*.json},
\texttt{results/oft/libero\_pro/*.json}, and
\texttt{results/oft/parity/*.json}. Wilson 95\% confidence intervals
and bootstrap p-values are in \texttt{paper/numbers.json}.

\end{document}
